\subsection{Suffix Automaton: la estructura estrella}\label{rec:sam}

Resuelve casi todo sobre substrings en O(n). Usa \codigo{a}--\codigo{z};
las mayúsculas se convierten solas. Se construye \textbf{una vez}:
\codigo{SuffixAutomaton sam(s);}.

\tablafunc{
\codigo{SuffixAutomaton(s)} & construye el autómata & $O(n)$ \\
\codigo{distinct_substrings()} & substrings distintas & $O(n)$ \\
\codigo{count_occurrences(p)} & veces que aparece \codigo{p} & $O(|p|)$ \\
\codigo{longest_common_substring(t)} & substring común más larga & $O(|t|)$ \\
\codigo{kth_distinct(k)} & k-ésima substring lexicográfica & $O(26|ans|)$ \\
\codigo{distinct_with_min_occ(k)} & substrings con al menos \codigo{k} apariciones & $O(n)$ \\
}

\lstinputlisting{codigo/suffix_automaton.cpp}

\ojo{Memoria O($26n$). \codigo{kth_distinct} devuelve \codigo{""} si
\codigo{k} no existe (vale hasta \codigo{10^18}). No agregues caracteres
después de construir sin recalcular \codigo{occ}.}